Here $\mathrm{MA}^{\mathrm{ext}}_\theta$ denotes the extended non-pluripolar Monge--Amp\`ere measure of \cref{thm:arbitrary-extended-ma}.

\begin{theorem}\label{thm:divisorial-residual-equation}\index{Divisor!residual equation}
Suppose that $\beta_D$ has a smooth positive representative, and put $c=\mathcal V_\mu(\theta;D)$. There is a $\theta$-plurisubharmonic function $u_0=\sigma_D+v$, with $v$ smooth, satisfying
\[
 \mathrm{MA}^{\mathrm{ext}}_\theta(u_0)=c\mu.
\]
Every $u\in\PSH(X,\theta)$ such that
\[
 \nu(u,D_j)\geq a_j\quad(1\leq j\leq N),\qquad
 \mathrm{MA}^{\mathrm{ext}}_\theta(u)\geq c\mu
\]
differs from $u_0$ by a constant. In particular, a solution of a normalized equation $\mathrm{MA}^{\mathrm{ext}}_\theta(u)=c'\mu$ in the ordinary singularity type of $\sigma_D$ has scalar $c'=c$, has a smooth remainder $u-\sigma_D$, and is unique up to constants.
\end{theorem}
\begin{proof}
\emph{Existence.} The smooth equation and \cref{thm:gaud-volume-dual-formula} give a smooth $v$ with
\[
 \Omega=\beta_D+\ddc v>0,\qquad \Omega^n=c\mu.
\]
Then $\theta+\ddc(\sigma_D+v)=[D]+\Omega\geq0$. Off $D$, the extended measure of $u_0=\sigma_D+v$ is $\Omega^n$, and both measures give zero mass to $D$. This proves existence.

\emph{Reduction to the residual background.} Let $u$ satisfy the two inequalities. Removing the divisor as in the proof of \cref{prop:divisorial-gaud-feasibility}, we obtain
\[
 w=u-\sigma_D-v\in\PSH(X,\Omega),\qquad
 \Omega+\ddc w=\theta+\ddc u-[D]\geq0,
\]
where the difference is first taken in $L^1$ and then replaced by its canonical representative. On $X\setminus D$, the two currents differ only by smooth changes in the decomposition of their local potentials, so their extended measures agree there, and both measures vanish on $D$. Hence
\[
 \mathrm{MA}^{\mathrm{ext}}_\Omega(w)
       =\mathrm{MA}^{\mathrm{ext}}_\theta(u)\geq c\mu.
\]

\emph{Rigidity.} Take $\mathrm{d}V=\mu$. The zero function solves the smooth equation $\Omega^n=c\mu$ in the background $\Omega$, so its subsolution threshold is $b_\Omega(1)=c$, by \cref{prop:individual-normalizer-lp}. For $n\geq2$, the normalized $(n-1)$st power of the Gauduchon metric conformal to $\Omega$ is a smooth strictly positive maximizer of the dual problem. By \cref{cor:smooth-dual-exact-equation}, any two threshold subsolutions differ by a constant; comparing $w$ with zero shows that $w$ is constant. For $n=1$, the positive current $\Omega+\ddc w$ dominates its extended measure, hence dominates $c\mu=\Omega$, while both have total mass $\int_X\Omega=c$ by Stokes' theorem. Thus $\ddc w=0$, and $w$ is again constant.

\emph{Solutions in the ordinary type.} Finally, let $u$ be in the ordinary singularity type of $\sigma_D$ with $\mathrm{MA}^{\mathrm{ext}}_\theta(u)=c'\mu$ for some $c'\geq0$. Then $\nu(u,D_j)=a_j$, and the same reduction gives a bounded $w\in\PSH(X,\Omega)$ with $\mathrm{MA}_\Omega(w)=c'\mu$. Suppose $c'<c$. For every constant $C$, the bounded functions $w$ and $C$ satisfy $\mathrm{MA}_\Omega(w)=(c'/c)\mathrm{MA}_\Omega(C)$, so \cref{thm:bounded-strict-domination} gives $w\geq C$, which is impossible for a bounded $w$. The case $c'>c$ is symmetric, with the roles of $w$ and the constants exchanged. Therefore $c'=c$, and the rigidity just proved gives $u=u_0+\mathrm{constant}$.
\end{proof}
